\section{Consul of God: Governance, Justice, War, and Empire (590--604)}

\subsection{The \emph{Registrum} and the patrimony of the poor}

More than eight hundred fifty letters survive from Gregory's fourteen years as
pope, arranged by his later Benedictine editors into fourteen books matching
the fourteen Roman indiction years of his reign (beginning each September from
Indiction IX in September 590 to Indiction VII in 603--604; Barmby,
Prolegomena).\footnote{Modern critical editions (beginning with Ewald and
Hartmann in the MGH) reorder several epistles within the books, so epistle
numbers in the Benedictine/Migne arrangement used by Barmby occasionally differ
by one or two digits from MGH numbering (for example, Benedictine \work{Ep.}\
I.35, I.43, and I.47 correspond to MGH I.34, I.41, and I.45). Citations here
follow the Benedictine numbering printed in NPNF2-12.} Through this
correspondence, Gregory stands before us in unguarded daily detail: checking
grain weights, rebuking bishops, comforting bereaved noblewomen, disciplining
wayward monks, negotiating truces, and funding hospitals.

Nowhere is his administrative genius clearer than in his oversight of the
\latin{Patrimonium Petri}---the scattered estates endowed to the Roman Church
across Sicily, Campania, Calabria, Dalmatia, and southern Gaul. Their income
and harvests fed Rome's swollen refugee population, supported clergy and
monasteries, and kept famine at bay. In a remarkable letter of May 591 to
Subdeacon Peter, his rector in Sicily (\work{Ep.}\ I.44), Gregory issued a
sweeping charter of agrarian justice for the Church's tenant peasants
(\latin{rustici} or \latin{coloni}):

\begin{itemize}
\item \textbf{True weights and fair prices:} Peter was ordered to smash
oversized measuring vessels, cap the grain \latin{modius} at eighteen
\latin{sextarii}, abolish unlawful surcharges, and adjust cash commutations of
the corn tithe to the actual market price in times of plenty.
\item \textbf{Protection from usurers:} Because peasants had to pay the
\latin{burdatio} land-tax before selling their harvest and were being ruined by
pawnbrokers, Peter was instructed to advance the tax money from the Church's
own funds and let the peasants repay gradually after harvest.
\item \textbf{Household rights and restitution:} Peasant marriage fees were
capped at one gold \latin{solidus} (and less for the poor); relatives of
deceased tenants were guaranteed inheritance of their goods rather than seeing
estates confiscated by the Church; and anything extorted from a peasant by an
estate manager was to be restored directly to the victim, ``lest we ourselves
should seem to be abettors of violence,'' for ``we would not have the treasury
of the Church defiled by base gains.''
\item \textbf{Written charters of security:} Gregory ordered Peter to give the
peasants signed schedules of their exact legal dues and have his papal
instructions read aloud across every estate ``that they may know in what
points to defend themselves, under our authority, against acts of wrong,''
reminding Peter of ``the terrible Judge who is coming.''
\end{itemize}

The same letter displays Gregory's pastoral realism. Three years earlier,
married Sicilian subdeacons had been suddenly ordered to observe Roman
continence. Gregory judged it ``hard and improper'' to force separation on men
who had not promised chastity before ordination; he forbade any future
ordination without a prior promise, praised those who had already lived
continently, and left previously married men undisturbed while withholding
higher promotion (\work{Ep.}\ I.44). Elsewhere he routinely authorized selling
even consecrated liturgical silver to ransom prisoners of war, writing to
Bishop Donus of Messina that it is wrong ``for an exceedingly desolated Church
to prefer its property to its captives, or to loiter in redeeming them''
(\work{Ep.}\ VII.38).

\subsection{Siege, diplomacy, and imperial friction}

War quickly turned the bishop of Rome into the city's de facto civil defender.
By late 591 and 592, Duke Ariulf of Spoleto was marching toward Rome while the
imperial Exarch at Ravenna, Romanus, withdrew troops from Rome and left the
remaining Theodosian garrison unpaid. Writing to John of Ravenna (\work{Ep.}\
II.46), Gregory described his grief falling into physical sickness: Romanus
``both evades fighting against our enemies and also forbids our making peace.''
Acting on his own initiative, Gregory arranged a truce with Ariulf to save
Rome. Soon afterward, however, Romanus marched south and seized cities along
the Perugia corridor, provoking the Lombard King Agilulf to descend from Pavia
and lay siege to Rome itself (593).

When Emperor Maurice later rebuked Gregory as ``simple''---politely calling him
silly---for having been deceived by Ariulf and faulted Rome's defenders over
grain shortages, Gregory answered the emperor in June 595 with fearless
dignity (\work{Ep.}\ V.40). He recounted what he had witnessed during Agilulf's
siege: ``I saw with my own eyes Romans tied by the neck with ropes like dogs,
to be taken to France for sale.'' Defending the urban prefect Gregory and the
military commander Castorius, who had guarded the walls tirelessly beside him,
he told the emperor plainly: ``Yet this I briefly say, that, unworthy sinner as
I am, I rely more on the mercy of Jesus when He comes than on the justice of
your Piety. And there are many things that men are ignorant of with regard to
this judgment; for perhaps He will blame what you praise, and praise what you
blame'' (\work{Ep.}\ V.40).

To Empress Constantina in the same month (\work{Ep.}\ V.41), he exposed the
rapacity of imperial tax collectors in Sardinia, Sicily, and Corsica---where
exactions drove parents to sell their own children and flee to the Lombards.
Even if those taxes paid for Italy's defense, Gregory insisted, ``it is better
that we should not live temporally, than that you should find any hindrance in
the way of eternal life.'' To Bishop Sebastian of Sirmium he wrote bluntly that
the Exarch Romanus's malice ``has surpassed the swords of the Lombards''
(\work{Ep.}\ V.42). When Maurice issued a law forbidding soldiers and civil
officials to enter monasteries, Gregory displayed his exact view of church and
empire: as a loyal subject he distributed the emperor's decree, yet as a
watchman for God he wrote boldly to Maurice (\work{Ep.}\ III.65) warning that
Christ---who had raised Maurice from notary to emperor---would demand an
account if his law barred men from serving the heavenly King. After five years,
a compromise was reached (\work{Ep.}\ VIII.5, reported by Barmby). Meanwhile,
Gregory patiently cultivated the Catholic Bavarian princess Theodelinda, widow
of King Authari and wife of Agilulf. Through her mediation, a general Lombard
truce was finally concluded in 598--599 (to March 601) and renewed in September
603 (Barmby, Prolegomena, citing \work{Ep.}\ IX.42--43, X.37, XIV.12 and Paul
the Deacon, \work{Historia Langobardorum} IV.8, 29).

\subsection{Jews, schismatics, and the limits of sixth-century tolerance}

Gregory's governance must be read inside the legal and religious horizon of the
Christian Roman Empire, without either modernizing him or flattening his real
distinctions. His policy toward Jewish communities was notably principled and
became a cornerstone of medieval papal protection. When Bishop Peter of
Terracina twice expelled local Jews from their synagogue, Gregory ordered them
restored and laid down a clear rule: ``as they are permitted, in accordance
with justice, to live under the protection of the Roman laws, let them keep
their observances as they have learnt them, no one hindering them''
(\work{Ep.}\ I.10), adding in a follow-up letter that ``those who dissent from
the Christian religion must needs be gathered together to unity of faith by
gentleness, kindness, admonition, persuasion'' (\work{Ep.}\ I.35). Hearing in
June 591 that Jews travelling to Marseilles were being coerced into baptism, he
wrote immediately to Bishops Virgilius of Arles and Theodorus of Marseilles to
forbid forced conversion:

\begin{quote}
For, when any one is brought to the font of baptism, not by the sweetness of
preaching, but by compulsion, he returns to his former superstition, and dies
the worse from having been born again. Let, therefore, your Fraternity stir up
such men by frequent preaching, to the end that through the sweetness of their
teacher they may desire the more to change their old life. (\work{Ep.}\ I.47)
\end{quote}

At the same time, Gregory strictly enforced imperial Roman law prohibiting Jews
from purchasing or retaining Christian slaves---freeing any slave who declared
a desire for Christianity without compensation to the owner (\work{Ep.}\ I.10;
Barmby, Prolegomena, citing III.38, IV.9, 21, VI.32, VII.24, IX.36, 110)---and
approved lowering rents for Jewish tenants on Church lands who embraced
Christianity (Barmby, citing II.32, V.8). Toward Christian schismatics and
rural pagans, moreover, his approach was far sterner. He repeatedly urged
imperial authorities to enforce existing penal laws against the lingering
Donatist schism in North Africa (\work{Ep.}\ I.74), and on Church estates in
Sardinia and Sicily he authorized heavy fiscal pressure against obstinate pagan
peasants and imprisonment or corporal punishment for idolaters and diviners
(Barmby, Prolegomena, citing IV.26 and IX.65). Even in the Three Chapters
controversy---where Istrian and Aquileian bishops broke communion with Rome
over the Fifth Council---Gregory initially sought imperial compulsion before
Maurice ordered him to desist (\work{Ep.}\ I.16, II.46), and despite his warm
friendship with Queen Theodelinda, he died without convincing either her or
Aquileia to accept the council (\work{Ep.}\ XIV.12, reported by Barmby).

\subsection{Servant of the servants of God: the ``Universal Bishop'' controversy}

No ecclesiastical struggle absorbed more of Gregory's energy than the dispute
over the title \latin{Oikoumenikos} (``Ecumenical'' or ``Universal''
Patriarch) adopted by Patriarch John IV the Faster of Constantinople at a synod
in 587 and retained by his successor Cyriacus (596--606). Pope Pelagius II had
already annulled the 587 synod's acts; when John sent Gregory court records in
which the title appeared ``almost in every line'' (\work{Ep.}\ V.19), Gregory
launched a sustained campaign in 595 (dated to the Kalends of January 595 in
Barmby's Benedictine text of \work{Ep.}\ V.18, and placed in June 595 in the
MGH chronological arrangement) through five decisive letters to John
(\work{Ep.}\ V.18), his envoy Sabinianus (V.19), Emperor Maurice (V.20),
Empress Constantina (V.21), and Patriarchs Eulogius of Alexandria and
Anastasius of Antioch (V.43).

Gregory's argument joins two convictions that later polemicists have often
pulled apart. On one hand, he asserted the divinely given Petrine primacy of
the Apostolic See with unmistakable clarity. Writing to Maurice, he cited the
three Petrine texts of the Gospels---John 21:17 (``Feed My sheep''), Luke 22:31
(``strengthen thy brethren''), and Matthew 16:18--19 (the Rock and the
keys)---declaring that ``the care and principality of the whole Church is
committed to him'' (\work{Ep.}\ V.20). He entertained appeals from Eastern
presbyters against Constantinople (\work{Ep.}\ VI.14--17) and insisted that
Constantinople itself was subject to the Apostolic See (Barmby, citing
\work{Ep.}\ IX.12). Yet he also described Rome, Alexandria (founded by Peter's
disciple Mark), and Antioch (where Peter first sat) as sharing one Petrine
chair: ``it is the See of one, and one See, over which by Divine authority
three bishops now preside'' (\work{Ep.}\ VII.40).

On the other hand, Gregory saw in the title ``Universal Bishop'' an attack on
the dignity of every other bishop in the Church---as though one prelate were
the sole bishop and all others merely his deputies. Even Peter, he reminded
Maurice, ``received the keys of the heavenly kingdom \dots\ and yet he is not
called the universal apostle'' (\work{Ep.}\ V.20). Comparing the proud title to
Lucifer's boast in Isaiah 14:13, he warned Maurice in 597 that ``whosoever
calls himself, or desires to be called, Universal Priest, is in his elation the
precursor of Antichrist'' (\work{Ep.}\ VII.33). Crucially, Gregory refused the
same honor for himself. When Eulogius of Alexandria addressed him as
``Universal Pope'' and wrote ``as you have commanded,'' Gregory rebuked him
gently and memorably:

\begin{quote}
This word, command, I beg you to remove from my hearing, since I know who I am,
and who you are. For in position you are my brethren, in character my fathers.
\dots\ Nor do I regard that as an honour whereby I know that my brethren lose
their honour. For my honour is the honour of the universal Church: my honour is
the solid vigour of my brethren. Then am I truly honoured when the honour due
to all and each is not denied them. For if your Holiness calls me Universal
Pope, you deny that you are yourself what you call me universally. But far be
this from us. Away with words that inflate vanity and wound charity.
(\work{Ep.}\ VIII.30)
\end{quote}

In place of inflated titles, Gregory popularized the papal style that his
successors still bear: \latin{servus servorum Dei}, ``servant of the servants
of God'' (attested in the salutations of \work{Moralia}, Ep.\ Ded., \work{Ep.}\
I.36, VI.51, and the Preface to \work{Homiliae in Evangelia}). Historically,
his campaign did not persuade Constantinople to abandon \latin{Oikoumenikos},
nor did Eulogius and Anastasius break communion with the imperial city over a
Greek honorific that carried far less exclusive force in Greek usage than
\latin{universalis} sounded in Latin ears.

\subsection{The shadow of 603: the letters to Phocas}

The gravest stain on Gregory's epistolary record fell near the end of his
reign. In November 602, a mutiny on the Danube overthrew Emperor Maurice and
placed the brutal centurion Phocas on the throne; Maurice's five sons were
slaughtered before his eyes before the emperor himself was executed. In the
spring and summer of 603, Gregory sent effusive letters of congratulation to
Phocas and his wife Leontia (\work{Ep.}\ XIII.31, 38, 39, summarized at length
in Barmby's Prolegomena).

Barmby notes two mitigating historical facts without excusing the tone. First,
Gregory had no resident \latin{apocrisiarius} in Constantinople at the close of
Maurice's reign (\work{Ep.}\ XIII.38), so it remains uncertain how fully he
knew the atrocities of the coup when he wrote. Second, late Roman diplomatic
protocol required bishops to greet a newly enthroned Augustus in high panegyric
formulas---just as Gregory had earlier used courtly flattery toward the
ruthless Merovingian queen Brunehild (\work{Ep.}\ VI.5, 59) to win support for
the Church in Gaul and the English mission. Frustration over Maurice's refusal
to aid Rome against the Lombards or restrain Cyriacus, combined with the
torment of advanced gout, may also have embittered Gregory's memory of the dead
emperor. Even with every allowance made, however, Gregory's jubilant welcome of
Phocas contrasts painfully with his earlier praise of Maurice as ``most
Christian of princes'' (\work{Ep.}\ VI.16) and remains a sobering reminder that
sanctity and statesmanship do not immunize a pastor against grave political
misjudgment.
